\scanpage{5}
Descriptions conjecturales des motifs sur corps de nbs \deleted{$\unsure{E}$} $k$

1) Double \unsure{viseur|vision} Betti--De Rham

2) Double \unsure{viseur|vision} Betti--Hodge

3) \deleted{Betti + \unsure{bigrad}} structure de Hodge $(k=\overline{\mathbf Q})$

4) \deleted{Betti + \unsure{connexion}, \unsure{viseur|vision} Betti}
\inserted{Betti -- (Gauss)Manin} $(k=\overline{\mathbf Q})$

5) \deleted{Tate} Betti--Tate

6) \deleted{De Rham--cristal} De Rham filtré cristallin

Description conjecturale des groupes $p$-divisibles \inserted{mod isogénie}
(plus gén$^t$ des pseudo-motifs) sur un corps local
[\deleted{\unsure{spectre|partie}}\inserted{corps de fr.} d’un ann. de val. discrète complet
à corps rés \unsure{$p$-fait}] :

3') structure de Hodge--\inserted{Tate ; module $N$ de t.f. sur $\mathbf Q_p$,
et bigr.d. sur $N\otimes_{\unsure{\mathbf Q_p|\mathbf Q}}\mathbf C$}
(mod ext. finie de $K$)

2') double \unsure{viseur|vision} Tate--Hodge

5') Tate (module galoisien $\ell$-adique)

6') \deleted{module \unsure{vect. de di}} cristal filtré
[vectoriel \deleted{\unsure{sur}} dim finie sur $K$, une \unsure{application|bijection}
$\sigma$-semilinéaire, et filtration, graduation]

Par quoi remplacer $1',4'$ ??
\margin{4}
